\section{Experimental Setup}

\subsection{Retrieval and Evaluation Protocol}
For all datasets, we use the PubMed Entrez API to execute the generated Boolean queries and retrieve candidate documents by matching on PMIDs~\cite{sayers2010general}. This simulates a realistic literature search workflow, allows us to evaluate query effectiveness in a practical retrieval setting, and follows standard methodology from previous work~\cite{wang2023can, wang2025reassessing}.
We apply the same query validity check protocol described in Section~\ref{sec:retrieval_reward} to detect and reject malformed queries. Queries that fail this check are considered invalid, and the model is prompted to regenerate until a valid query is produced, with a maximum of 10 attempts.%
\footnote{We cap regenerations at 10 to avoid excessive API usage and model inference.}
This validation process aligns with established evaluation practices~\cite{wang2025reassessing}.

We evaluate with measures used in earlier research~\cite{wang2025reassessing}. The \textbf{primary evaluation metrics} are recall, F\textsubscript{3}, and the percentage of queries obtaining recall above 80\% and 90\%. These reflect the high-recall requirements of systematic reviews, where omissions of relevant studies can critically undermine review quality.
As \textbf{secondary metrics}, we report precision, the average number of documents retrieved (to measure screening effort), the average number of regenerations per query, and the success rate under 10 attempts---to capture robustness and generation stability beyond recall-focused evaluation.

\subsection{Model Variants}
Our primary experiments are conducted using models from the Qwen3 family~\cite{yang2025qwen3}. Unless otherwise specified, all trained AutoBool Models are based on \texttt{Qwen3-4B}, which serves as the default backbone throughout our main results. These models are fine-tuned using our reward-driven training framework (Section~\ref{sec:retrieval_reward}) to optimize for systematic review retrieval effectiveness. GRPO is applied to guide query generation behavior based on retrieval performance signals. 
At inference time, we use the same prompt as during training and fix the decoding temperature to 0.6, following the Qwen3 recommendations~\cite{yang2025qwen3}.


\subsection{Evaluation Datasets}
We evaluate our models on three datasets: the \textbf{PubTemp set}, and two established benchmarks: the \textbf{CLEF TAR} collection~\cite{kanoulas2018clef} and the \textbf{Seed Collection}~\cite{wang2022little}.

We follow prior work~\cite{wang2025reassessing} using the CLEF TAR 2017 and 2018 subsets (72 topics total) and the Seed Collection (40 topics). Each topic is defined by experts and paired with a manually curated set of relevant studies. These benchmarks have been widely adopted for evaluating automatic Boolean query generation in the context of systematic review automation~\cite{macfarlane2022search, kusa2023csmed, wang2022neural, stevenson2023stopping, lee2018seed}.

\subsection{Training Parameters}

All models are trained using the GRPO RL algorithm with a retrieval-based reward. We adopt LoRA-based parameter-efficient fine-tuning and use the vLLM backend in colocated mode.
For each prompt the model generates four completions to compute reward scores. We use an effective batch size of 16 via gradient accumulation. Prompt and completion length limits are 768/1024 tokens for non-reasoning prompts and 1024/3072 for reasoning-based prompts.
Unless otherwise noted, we set the training temperature to 1.2; its effect is analyzed in Section~\ref{sec:ablation}. Full training hyperparameters are listed in Table~\ref{appendix:tab-training-details} in Appendix~\ref{appendix:training-details}.